\section{Method}
\label{sec:method}

This section defines every quantity used later in the paper: extraction, the
vocabulary, the components and scores, the reasons for three design choices,
and panel inference (\S\ref{sec:method-extraction}--\ref{sec:method-inference}).

\subsection{Extraction of ESG passages}
\label{sec:method-extraction}

The reports of \S\ref{sec:data} are annual reports and universal registration
documents, in which ESG disclosure is embedded in a much larger financial and
legal document. Scored whole, such a document would let financial-statement
bulk and document length dominate every component. We therefore first extract
its ESG passages; every later measure is computed on them alone.

Each PDF is read as text blocks in reading order, and a block becomes a
paragraph only after five filters: blocks entirely within the top or bottom
8\,\% of the page are dropped; short margin blocks that recur, digits aside, on
more than 30\,\% of a document's pages are dropped as running headers or
footers; tabular blocks (a high share of digits in a few lines, or short
lines of which the first lacks terminal punctuation) are dropped unless they
contain a vocabulary match; short all-capital blocks are taken as section
headings, which label the paragraphs that follow instead of entering the text;
and navigation text (tables of contents, GRI and
ESRS correspondence tables, stray page references), recognised by dotted
leaders, trailing page numbers, multi-level section numbers and disclosure
codes, is dropped, because such lists of ESG section names would otherwise
pass the density rule below. Paragraphs are then stripped of web addresses,
rejoined across line-break hyphenation and collapsed to single spaces.

Each paragraph is matched against the vocabulary of \S\ref{sec:method-vocab}
on its unlemmatised text, ignoring case, with spaces and hyphens
interchangeable between the words of a phrase and longer entries matched before
the shorter ones they contain, so that \emph{carbon footprint} counts once. A
paragraph with at least one match per 100 words is retained together with one
paragraph on each side; overlapping or adjacent windows are merged into a
single zone, and zones shorter than 150 characters are discarded. The zones of
a report, in document order, form its extracted text. Across the 343 reports,
extraction kept 66,171 zones holding 28.1 million of the 71.0 million words of
the PDF text layers, 39.6\,\% (per report, median 38.3\,\%, range
7.1--75.6\,\%). The procedure depends only on the PDF and the vocabulary, and
its recall is bounded by the vocabulary: a passage that uses none of its
entries is never selected.

\subsection{ESG vocabulary}
\label{sec:method-vocab}

Candidate entries came from three sources: the disclosure topics and
terminology of reporting frameworks and standards (GRI, SASB, TCFD and the
ESRS), including the names of the frameworks and regulatory regimes under
which firms label their disclosure; the ESG vocabulary of dictionary-based
studies of corporate disclosure; and a structured review of sample reports from
the corpus, in which independent readers listed the ESG terms they found and
independent evaluators judged the pooled list.
\textbf{[Author: identify the reviewers of the candidate lists.]}
A candidate was admitted if its occurrences in the extracted text, read in
context, predominantly carry an ESG sense (a disclosure, target, data point or
policy on an environmental, social or governance matter). Candidates from the
review had in addition to be dispersed across firms, with at least 50
occurrences in the reports of at least 5 of the 49 firms. A polysemous head was
admitted only inside enumerated collocations (\emph{carbon footprint} rather
than \emph{footprint}; the ESG collocations of \emph{social} and \emph{water}),
and a morphological family or a term carrying digits was written as a pattern
(every word beginning \emph{sustainab-} or \emph{recycl-}; CO$_2$e; Scope 1--3;
ISO standard numbers). Each entry belongs to one pillar: environmental (E),
social (S), governance (G), or transversal (TRANS) for umbrella terms and
reporting regimes (\emph{sustainability}, \emph{ESG}, \emph{CSRD}). Entries
whose occurrences concentrate in one industry are designated sectoral; they are
part of the vocabulary, and excluding them is a sensitivity check
(\S\ref{sec:robustness}). Table~\ref{tab:method-vocab} gives the composition.

Terms were admitted after reviewing a sample of their occurrences in context,
and the frequency-weighted precision of the resulting vocabulary is estimated
from those samples, not measured directly on the final extraction, at about
82\,\%. The stability of the scores under perturbations of the vocabulary is
reported in \S\ref{sec:robustness}.

\begin{table}[tbp]
\centering
\small
\caption{Composition of the ESG vocabulary. Terms and collocations are literal
phrases; a collocation stands in for a polysemous head. Sectoral entries cut
across the three type columns. Scoring entries are the distinct lemmatised
terms (patterns act at extraction only); the last column is their share of the
1,047,951 vocabulary occurrences in the extracted text of the 343 reports.}
\label{tab:method-vocab}
\begin{tabular}{@{}lrrrrrrr@{}}
\toprule
 & \multicolumn{5}{c}{All entries} & \multicolumn{2}{c}{Scoring entries} \\
\cmidrule(lr){2-6}\cmidrule(l){7-8}
Pillar & Terms & Colloc. & Patterns & Total & Sectoral & Entries & Occ.\ (\%) \\
\midrule
E (environmental)   & 123 & 29 & 16 & 168 & 14 & 152 & 38.2 \\
S (social)          & 105 & 23 & 11 & 139 & 14 & 127 & 17.6 \\
G (governance)      &  65 & 27 &  3 &  95 &  8 &  90 & 30.0 \\
TRANS (transversal) &  27 &  3 &  2 &  32 &  0 &  30 & 14.2 \\
\midrule
Total               & 320 & 82 & 32 & 434 & 36 & 399 & 100.0 \\
\bottomrule
\end{tabular}
\end{table}

One vocabulary drives both extraction and scoring: all 434 entries select the
passages, and the 402 terms, lemmatised, are what the substance component
counts. Changing the vocabulary therefore changes the text on which every
component is computed, not only the substance count. For scoring, the extracted
text is lowercased, tokenised and lemmatised with a general-purpose English
language model; stopwords (the standard list and 189 corpus-specific words),
punctuation, numerals and tokens shorter than three characters are removed, and
the vocabulary passes through the same processing, so that each entry is
matched in the representation of the text it is counted on. Eleven phrases that
this processing would empty or reduce to a generic word (\emph{well-being},
both of whose words are stopwords; \emph{say on pay} to \emph{pay};
\emph{ISO 14001} to \emph{iso}) are first replaced by single protected tokens,
in text and vocabulary alike. The 402 terms yield $V = 399$ distinct
lemmatised scoring entries.

\subsection{Components and scores}
\label{sec:method-index}

Let $d = 1, \dots, n$ index the $n = 343$ reports and $j = 1, \dots, V$ the
scoring entries. Each report has two representations of its extracted text:
the \emph{lemmatised text} of \S\ref{sec:method-vocab}, with $L_d$ tokens, and
the \emph{raw text}, with case, digits and stopwords intact, with $R_d$
tokens. Let $c_{dj}$ be the number of occurrences of entry $j$ in the
lemmatised text of $d$, counted over word sequences of one to four tokens (the
length of the longest entry). Entries are counted independently, so an
occurrence of a multi-word entry also counts for any shorter entry it contains.

\paragraph{Substance.} $\SUS$ is the density of ESG mentions, the number of
vocabulary occurrences per 100 lemmatised tokens, with uniform term weights:
\[
\SUS_d \;=\; \susdens_d \;=\; 100 \cdot \frac{\sum_{j} c_{dj}}{L_d},
\]
written $\susdens$ where it is compared with the alternatives of
\S\ref{sec:method-choices}. It uses no document frequencies, so a report's
value does not depend on the other reports in the corpus.

\paragraph{Tone.} $\SEN$ is the Loughran--McDonald polarity
\citep{loughran2011} of the lemmatised text,
\[
\SEN_d \;=\; \frac{s^{+}_d - s^{-}_d}{s^{+}_d + s^{-}_d + \varepsilon},
\qquad \varepsilon = 10^{-6},
\]
where $s^{+}_d$ and $s^{-}_d$ count the tokens whose stem is in the
dictionary's positive and negative lists, and $\varepsilon$ only prevents
division by zero. It is a bounded ratio on $[-1, 1]$ that depends on the
balance of positive and negative words, not on how many there are.

\paragraph{Quantification.} $\QUANT$ is the number of pattern matches per raw
token, summed over four families: percentages; numbers of four or more digits
other than years; physical units (tonnes, CO$_2$e, GHG, energy, power and
volume units); and the names of reporting frameworks (GRI, TCFD,
SASB, CSRD, SDG, the taxonomy and others). It is computed on the raw text
because preprocessing removes numerals.

\paragraph{Hedging.} $\HEDGE$ is the share of the words of the raw text
(letter runs, hyphenated compounds kept whole) that belong to the Uncertainty,
Weak Modal, Strong Modal and Constraining lists of the Loughran--McDonald master
dictionary, 495 word forms taken as written. It is computed on the raw text
because the modal verbs (\emph{may}, \emph{might}, \emph{could}, \emph{must},
\emph{will}) are stopwords. Matching is by form, not by sense: the five most
frequent forms (\emph{risk}, \emph{risks}, \emph{will}, \emph{may},
\emph{requirements}) make up 49.4\,\% of all matches, so $\HEDGE$ registers to
a large extent how often a report names risk.

\paragraph{Breadth and pillars.} $\BREADTH_d = \exp\bigl(-\sum_j p_{dj}\ln
p_{dj}\bigr)$, with $p_{dj} = c_{dj} / \sum_k c_{dk}$, is the effective number
of distinct entries a report uses, a Hill number of order one
\citep{hill1973}; it is a diagnostic, not a component of either score. The
pillar density $100 \sum_{j \in p} c_{dj} / L_d$ counts the mentions of
pillar-$p$ entries per 100 lemmatised tokens, and the four pillar densities sum
to $\SUS_d$; for a set of reports, the pillar share pools the mentions of the
set. Both are descriptive and enter neither score.

\paragraph{Scores.} Each component $x$ is rescaled across the corpus by a
$z$-score, $\Z(x)_d = (x_d - \bar{x}) / \sigma_x$, with $\bar x$ the mean and
$\sigma_x$ the population standard deviation over all 343 reports pooled across
the seven years, so that yearly means of a score are comparable. The two
scores are
\begin{align}
\label{eq:esgsi}
\ESGSI_d    &= \Z(\SEN)_d - \Z(\SUS)_d, \\[2pt]
\label{eq:esgsiext}
\ESGSIext_d &= \Z(\SEN)_d - \Z(\SUS)_d - w_q\,\Z(\QUANT)_d + w_h\,\Z(\HEDGE)_d,
\qquad w_q = w_h = 0.5 .
\end{align}
$\ESGSI$~\eqref{eq:esgsi} is the main score of the paper. It keeps the
architecture of \citet{lagasio2024}, tone minus substance with each rescaled
across the sample, with substance measured by density and tone by a
finance-specific dictionary in place of a general-purpose polarity score. The
Extended ESGSI~\eqref{eq:esgsiext} is the new metric this paper proposes. Quantified disclosure is further evidence of
substance, so $\Z(\QUANT)$ is subtracted, on the same side as $\Z(\SUS)$;
hedging qualifies whatever a report asserts, so $\Z(\HEDGE)$ is added, on the
side of tone. The weights give each added component half the weight of an
original one, so that together they weigh as much as one of them and the
tone--substance contrast remains the core of the score. They are a stated
choice, not an estimate, since no external criterion is available to calibrate
them; \S\ref{sec:robustness} reports the scores over a grid of weights.

Both scores are continuous and corpus-relative: each has mean zero by
construction, and its sign says only whether a report lies above or below the
corpus mean of the combination. We compare positions within the corpus:
differences of means, trends, rank correlations and the overlap of the top and
bottom deciles.

\subsection{Why density, \texorpdfstring{$z$}{z}-scores and no threshold}
\label{sec:method-choices}

\citet{lagasio2024} computes substance with a TF-IDF vectoriser
\citep{pedregosa2011} over an author-built keyword list and discloses four of
its settings: two document-frequency bounds, a vocabulary cap and the range of
word sequences. It does not state whether each document's term vector is
scaled to unit Euclidean length (vector normalisation), nor how the per-term
values become a document score. The aggregation is immaterial once scores are
standardised, since a sum and a mean over the vocabulary differ by the factor
$1/V$. Vector normalisation is the vectoriser's default and all four disclosed
settings depart from defaults; inferring, as a reading and not a finding, that
only departures were reported, the closest reading of the antecedent is
$\suslag$ below. We also compute $\sustfidf$, which differs from $\susdens$
only in the idf weights:
\[
\suslag_d = \frac{1}{V}\sum_j \frac{c_{dj}\,\mathrm{idf}_j}
{\lVert \mathbf{c}_d \odot \mathbf{idf} \rVert_2},
\qquad
\sustfidf_d = 100 \cdot \frac{\sum_j c_{dj}\,\mathrm{idf}_j}{L_d},
\]
where $\mathrm{idf}_j = \ln[(1+n)/(1+\mathit{df}_j)] + 1$ and $\mathit{df}_j$
is the number of reports containing entry $j$.

A term vector and any positive multiple of it have the same unit vector, so
$\suslag$ keeps the mix of terms and discards their amount, a property suited
to retrieval \citep{sparckjones1972,manning2008} but not to a score that sets
tone against the amount of substance behind it. Panel~A of
Table~\ref{tab:method-choices} shows this on the report at the median of
$\susdens$: multiplying its ESG counts by up to 20 raises its density from 6.78
to 33.79 and leaves $\suslag$ unchanged, and diluting it with 152,378 non-ESG
tokens divides its density by 3.5 and leaves $\suslag$ identical. What remains
after vector normalisation is dispersion over terms: the sum of a non-negative
unit vector lies between 1 and $\sqrt{V}$ and grows as mentions spread over more
entries. On our corpus $\suslag$ correlates 0.955 with $\BREADTH$ and 0.405
with $\susdens$; under vector normalisation, the substance component measures
topical breadth, not the amount of disclosure.

\begin{table}[tbp]
\centering
\small
\caption{Why density. Panel~A: a synthetic document written out from the ESG
counts of the report at the median of $\susdens$, with every count multiplied
by $k$, and the real report with 152,378 non-ESG tokens appended; all 343
reports re-vectorised in every run; $\suslag$ relative to the first row of each
block (maximum absolute deviation $3.5\times10^{-18}$). Panel~B: Pearson
correlation with $\QUANT$ ($n = 343$, 49 firms); disjoint bases remove from the
vocabulary the 17 entries also matched by a $\QUANT$ pattern (399 to 382
entries) and restrict $\QUANT$ to percentages and large numbers; 95\,\%
percentile intervals from 9,999 bootstrap samples of firms, each with its
seven reports.}
\label{tab:method-choices}
\begin{tabular}{@{}lrrrr@{}}
\toprule
\multicolumn{5}{@{}l}{\emph{A. Amount of ESG content and the substance specifications}} \\
 & Tokens & ESG mentions & $\susdens$ & $\suslag$ (rel.) \\
\midrule
Amplified, $k = 1$  &  60,951 &  4,131 &  6.78 & 1.000000 \\
Amplified, $k = 2$  &  70,612 &  8,262 & 11.70 & 1.000000 \\
Amplified, $k = 5$  &  99,595 & 20,655 & 20.74 & 1.000000 \\
Amplified, $k = 20$ & 244,510 & 82,620 & 33.79 & 1.000000 \\
\addlinespace
Real report         &  60,951 &  3,984 &  6.54 & 1.000000 \\
Diluted             & 213,329 &  3,984 &  1.87 & 1.000000 \\
\midrule
\multicolumn{5}{@{}l}{\emph{B. Convergent validity: correlation with $\QUANT$}} \\
 & \multicolumn{2}{c}{Full bases} & \multicolumn{2}{c}{Disjoint bases} \\
\cmidrule(lr){2-3}\cmidrule(l){4-5}
 & $r$ & 95\,\% interval & $r$ & 95\,\% interval \\
\midrule
$\susdens$           & 0.6305 & $[0.523,\ 0.732]$ & 0.3992 & $[0.227,\ 0.539]$ \\
$\sustfidf$          & 0.6392 & $[0.550,\ 0.736]$ & 0.4007 & $[0.256,\ 0.535]$ \\
$\suslag$            & 0.2984 & $[0.041,\ 0.551]$ & 0.1850 & $[-0.072,\ 0.483]$ \\
$\susdens - \suslag$ & 0.3321 & $[0.075,\ 0.589]$ & 0.2142 & $[-0.143,\ 0.519]$ \\
\bottomrule
\end{tabular}
\end{table}

An empirical check needs a yardstick built by other means. $\QUANT$ counts
pattern matches on the raw text, with no vocabulary, vectoriser or document
frequencies, and a report that discloses substance should also quantify it
\citep{campbell1959}. On the full bases (Table~\ref{tab:method-choices},
panel~B) density correlates 0.63 with $\QUANT$ and $\suslag$ 0.30, and the
firm-bootstrap interval of the difference excludes zero. The two instruments
share 17 entries (framework acronyms, the taxonomy family, CO$_2$, GHG, ISO
standards) and read the same zones, so we repeat the test on disjoint bases.
Every correlation falls, by 37--38\,\%, and the ratio between the two
specifications does not (2.16 against 2.11), but the interval of the
difference now includes zero: convergent validity corroborates a density
specification without carrying the choice on its own, which rests on the
invariance above. The two density specifications are not separated by this
test (difference 0.009, interval $[-0.018, 0.044]$ on the full bases); we
adopt $\susdens$ because it estimates nothing on the sample.

The antecedent prints its index as a difference of $z$-scores but states that
the components are min--max rescaled, and its published statistics are
min--max: its mean index, $-0.2179$, is the difference of its component means,
$0.3727 - 0.5906$ \citep[p.~6]{lagasio2024}, whereas a difference of
$z$-scores has mean zero. Both score normalisations are increasing affine maps
and preserve ranks within a component; once two components are differenced,
they differ in how they equalise spreads and where they put zero. On our corpus
ranks barely move (Spearman correlation 0.997 between the $z$-score and
min--max versions of $\ESGSI$), but location does: under min--max the mean of
the score depends on where each component's mass sits within its observed
range ($+0.004$ with $\susdens$, $-0.203$ with $\suslag$, on the same
reports), and a single extreme report can shift it. A $z$-score puts each
component's zero at its corpus mean and measures it in its own standard
deviations, so a score's sign and size describe a position in the corpus and
nothing else. The antecedent also assigns every report to one of two classes by
the sign of its index. We apply no threshold of any kind. Zero is the corpus
mean of the combination, so a cut-off on a difference of standardised
components describes the relative location of two distributions, not the
conduct of a firm, and no external benchmark exists against which a boundary
could be calibrated for these reports (\S\ref{sec:limitations}).

\subsection{Inference on the panel}
\label{sec:method-inference}

Write $d = (i, t)$ for the report of firm $i = 1, \dots, M$ in year $t$, with
$M = 49$ in the balanced panel of \S\ref{sec:data}. Trends are estimated with
firm fixed effects, regressing the firm-demeaned outcome on the firm-demeaned
year; in the balanced panel the slope equals the mean of the firm-level
least-squares slopes. Standard errors are clustered by firm with the usual
small-sample correction, and $t$ statistics are referred to a $t$ distribution
with $M - 1$ degrees of freedom, $M$ being the number of firms in the sample
or subsample. Clustering allows any correlation among a firm's seven reports
but not a shock common to all firms in one year, which with seven years cannot
be clustered on; the slope summarises the yearly means. Because 49 clusters
are few, the fixed-effects slope is also tested with a restricted wild cluster
bootstrap \citep{cameron2008}: under a zero slope the residuals are the
demeaned outcomes, each of $B = 9{,}999$ replications multiplies all residuals
of a firm by one Rademacher weight ($\pm 1$ with equal probability), and the
$p$-value is the share of replications whose cluster-robust $t$ is at least as
large in absolute value as the observed one. Slopes by industry or country
interact the demeaned year with group indicators. Each equals the
mean firm slope of its group, but for a group of $M_k$ firms its clustered
variance rests on $M_k - 1$ free contributions and is too narrow for small
groups, so we read group slopes through the $t$ interval on the $M_k$ firm
slopes ($M_k - 1$ degrees of freedom). Equality across groups is tested with a
Wald test on the clustered variance, anti-conservative when a group has a
single firm. Yearly means carry 95\,\% intervals from a regression on year
indicators, clustered by firm.
